\specialchapter{INTRODUCTION TO CHAPTERS IV, V AND VI}

The study of semi-simple (analytic or algebraic) groups and their Lie algebras leads to the consideration of root systems, Coxeter groups and Tits systems. Chapters IV, V and VI are devoted to these structures.

To orient the reader, we give several examples below.

I. Let g be a complex semi-simple Lie algebra and h a Cartan subalgebra of \(g^{1}\) . A root of g with respect to h is a non-zero linear form \(\alpha\) on h such that there exists a non-zero element x of g with \([h, x] = \alpha(h)x\) for all \(h \in h\) . These roots form a reduced root system R in the vector space \(h^{*}\) dual to h. Giving R determines g up to isomorphism and every reduced root system is isomorphic to a root system obtained in this way. An automorphism of g leaving h stable defines an automorphism of \(h^{*}\) leaving R invariant, and every automorphism of R is obtained in this way. The Weyl group of R consists of all the automorphisms of \(h^{*}\) defined by the inner automorphisms of g leaving h stable; this is a Coxeter group.

Let \(G\) be a connected complex Lie group with Lie algebra \(\mathfrak{g}\) , and let \(\Gamma\) be the subgroup of \(G\) consisting of the elements \(h\) such that \(\exp_G(2\pi ih) = 1\) . Let \(R^-\) be the root system in \(\mathfrak{h}\) inverse to \(R\) , let \(Q(R^-)\) be the subgroup of \(\mathfrak{h}\) generated by \(R^-\) and let \(P(R^-)\) be the subgroup associated to the subgroup \(Q(R)\) of \(\mathfrak{h}^*\) generated by \(R\) (i.e.~the set of \(h \in \mathfrak{h}\) such that \(\lambda(h)\) is an integer for all \(\lambda \in Q(R)\) ). Then \(P(R^-) \supset \Gamma \supset Q(R^-)\) . Moreover, the centre of \(G\) is canonically isomorphic to \(P(R^-)/\Gamma\) and the fundamental group of \(G\) to \(\Gamma/Q(R^-)\) . In particular, \(\Gamma\) is equal to \(P(R^-)\) if \(G\) is the adjoint group and \(\Gamma\) is equal to \(Q(R^-)\) if \(G\) is simply-connected. Finally, the weights of the finite-dimensional linear representations of \(G\) are the elements of the subgroup of \(\mathfrak{h}^*\) associated to \(\Gamma\) .

\begin{enumerate}
\def\labelenumi{\Roman{enumi}.}
\setcounter{enumi}{1}
\tightlist
\item
  Let G be a semi-simple connected compact real Lie group, and let g be its Lie algebra. Let T be a maximal torus of G, with Lie algebra t, and let X be the group of characters of T. Let R be the set of non-zero elements \(\alpha\) of X such that there exists a non-zero element x of g with (Ad t). \(x = \alpha(t)x\) for all \(t \in T\) . Identify X with a lattice in the real vector space \(V = X \otimes_{Z} R\) ; then R is a reduced root system in V. Let N be the normaliser of T in G; the action of N on T defines an isomorphism from the group N/T to the Weyl group of R. We have \(P(R) \supset X \supset Q(R)\) ; moreover, \(X = P(R)\) if G is simply-connected and \(X = Q(R)\) if the centre of G reduces to the identity element.
\end{enumerate}

The complexified Lie algebra \(\mathfrak{g}_{(\mathbf{C})}\) of g is semi-simple and \(t_{(\mathbf{C})}\) is a Cartan subalgebra of it. There exists a canonical isomorphism from \(V_{(\mathbf{C})}\) to the dual of \(t_{(\mathbf{C})}\) that transforms R into the root system of \(\mathfrak{g}_{(\mathbf{C})}\) with respect to \(t_{(\mathbf{C})}\) .

\begin{enumerate}
\def\labelenumi{\Roman{enumi}.}
\setcounter{enumi}{2}
\tightlist
\item
  Let G be a connected semi-simple algebraic group over a commutative field k. Let T be a maximal element of the set of tori of G split over k and let X be the group of characters of T (the homomorphisms from T to the multiplicative group). We identify X with a lattice in the real vector space \(V = X \otimes_{Z} R\) . The roots of G with respect to T are the non-zero elements \(\alpha\) of X such that there exists a non-zero element x of the Lie algebra g of G with \((\operatorname{Ad} t)x = \alpha(t)x\) for all \(t \in T\) . This gives a root system R in V, which is not necessarily reduced. Let N be the normaliser and Z the centraliser of T in G and let \(N(k)\) and \(Z(k)\) be their groups of rational points over k. The action of \(N(k)\) on T defines an isomorphism from \(N(k)/Z(k)\) to the Weyl group of R.
\end{enumerate}

Let U be a maximal element of the set of unipotent subgroups of G, defined over k and normalised by Z. Put P = Z.U. Then \(P(k) = Z(k)\) . \(U(k)\) and \(P(k) \cap N(k) = Z(k)\) . Moreover, there exists a basis \((\alpha_{1}, \ldots, \alpha_{n})\) of R such that the weights of T in U are the positive roots of R for this basis; if S denotes the set of elements of \(N(k)/Z(k)\) that correspond, via the isomorphism defined above, to the symmetries \(s_{\alpha_{i}} \in W(R)\) associated to the roots \(\alpha_{i}\) , the quadruple \((G(k), P(k), N(k), S)\) is a Tits system.

\begin{enumerate}
\def\labelenumi{\Roman{enumi}.}
\setcounter{enumi}{3}
\tightlist
\item
  In the theory of semi-simple algebraic groups over a local field, Tits systems are encountered whose Weyl group is the affine Weyl group of a root system. For example, let \(\mathbf{G} = \mathbf{S}\mathbf{L}(n + 1, \mathbf{Q}_{p})\) (with \(n \geq 1\) ). Let B be the group of matrices \((a_{ij}) \in \mathbf{S}\mathbf{L}(n + 1, \mathbf{Z}_{p})\) such that \(a_{ij} \in pZ_{p}\) for i \textless{} j, and let N be the subgroup of G consisting of the matrices having only one non-zero element in each row and column. Then there exists a subset S of \(N/(B \cap N)\) such that the quadruple \((\mathbf{G}, \mathbf{B}, \mathbf{N}, \mathbf{S})\) is a Tits system. The group \(W = N/(B \cap N)\) is the affine Weyl group of a root system of type \(A_{n}\) ; this is an infinite Coxeter group.
\end{enumerate}

Numerous conversations with J. Tits have been of invaluable assistance to us in the preparation of these chapters. We thank him very cordially.
